\scanpage{7}
\uline{Corollaire} Ledit noyau est donc égal à l'image de $H^i(k,\operatorname{Ker}(F_{\unsure{G|C}}^n))$, $F_G^n:G-G^{(p^n)}$ le Frobenius itéré.

Or si $G$ est simple /$k$, alors $F_G^n$ est radiciel, donc son noyau est radiciel /$k$, a fortiori fini. Cela achève de prouver le Corollaire au lemme 1.

\uline{Lemme 3} Pour que $H^1(k,G)-H^1(k,G')$ soit un épim pour tt épim $G\xrightarrow{u}G'$ de schémas en groupes algébriques, il suffit qu'il en soit ainsi dans les deux cas suivants

\begin{enumerate}[label=\alph*)]
\item $G$, \deleted{$\unsure{K}$} et $K=\operatorname{Ker}u$ \deleted{\unsure{sont}} \inserted{(comme $G'$)} simples /$k$
\item $K$ est radiciel.
\end{enumerate}

Dans le cas général, il existe en effet un sous groupe radiciel \inserted{inv. $I$} de $G$ \unsure{[|(}noyau d'un Frobenius itéré convenable] tel que $G/I$ \deleted{$\unsure{K}$} et $K/I\cap K$ sont simples /$k$.

Considérons
\diagram{\begin{tikzcd}[ampersand replacement=\&]
e \arrow[r] \& K \arrow[r,dash] \arrow[d] \& G \arrow[r,dash,"\unsure{f|j}"] \arrow[d,"u"'] \& G/K \arrow[r] \arrow[d,"v"] \& e \\
e \arrow[r,dash] \& K/K\cap I \arrow[r,dash] \& G/I \arrow[r,dash,"f'"] \& G/KI \arrow[r,dash] \& e
\end{tikzcd}}
\deleted{Alors} d'où
% archive sheet number: 6

